\section*{الفصل \ref{ch:b2:amines} --- الأمينات ومركبات النيتروجين}

\begin{solution}{exo:b2:amines:1}
أولي؛ ثانوي؛ ثالثي؛ ملح أمونيوم رباعي؛ أولي (\omterm{def:b2:huckel:aromatic}{عطري}).
\end{solution}

\begin{solution}{exo:b2:amines:2}
\babelsublr{\foreignlanguage{arabic}{الإيثانأميد}~$<$~\foreignlanguage{arabic}{الأنيلين}~$<$~\foreignlanguage{arabic}{الأمونيا}~$<$~\foreignlanguage{arabic}{الميثيل أمين}} (قيم $\mathrm pK_a$ للأحماض المرافقة
$-0.6$، 4.60، 9.26، 10.66).
\end{solution}

\begin{solution}{exo:b2:amines:3}
\textit{N}-ميثيل هكسان حلقي \omterm{def:b2:amines:imine}{إيمين}، \ce{C6H10=N-CH3}: الأمين والكيتون في وسط ضعيف الحموضة
(pH 4--5)، مع إزالة الماء المتكوّن (عامل تجفيف أو تقطير أزيوتروبي).
\end{solution}

\begin{solution}{exo:b2:amines:4}
البروموبنزين؛ البنزونتريل؛ الفينول؛ 4-هيدروكسي آزوبنزين \ce{C6H5-N=N-C6H4-OH} (اقتران في الموضع بارا بالنسبة إلى
مجموعة \ce{O-} في الفينوكسيد).
\end{solution}

\begin{solution}{exo:b2:amines:5}
$\mathrm pK_a = 9.80$: عند pH 7 تكون النسبة $[\ce{(CH3)3NH+}]/[\ce{(CH3)3N}] = 10^{2.8} \approx 630$: فتسود
الصورة المبرتنة؛ وعصير الليمون (pH قرب 2) يبرتنه كله تقريبًا. ورجّ محلول
عضوي لأمين مع حمض مائي مخفف ينقل الأمين، على هيئة أيون أمونيومه،
إلى الماء؛ والقاعدة تحرره من جديد.
\end{solution}

\begin{solution}{exo:b2:amines:6}
1-(هكس-1-ين حلقي-1-يل)بيروليدين. بعد الإضافة وفقد الماء يحمل النيتروجين
شحنة موجبة (أيون إيمينيوم) ولا هيدروجين لديه ليفقده؛ فيُفقد بروتون من الكربون
المجاور بدلًا من ذلك، فتتكوّن الرابطة \ce{C=C}.
\end{solution}

\begin{solution}{exo:b2:amines:7}
البنزالدهيد مع البروبان-2-أمين، أو البروبانون مع البنزيل أمين، كل منهما مع \ce{NaBH3CN} عند pH
قرب 6.
\end{solution}

\begin{solution}{exo:b2:amines:8}
البروبيوفينون (1-فينيل بروبان-1-ون)، \ce{C6H5COCH2CH3}: إضافة واحدة تعطي أنيون \omterm{def:b2:amines:imine}{الإيمين}،
الذي يُحلمَه إلى الكيتون.
\end{solution}

\begin{solution}{exo:b2:amines:9}
فثاليميد البوتاسيوم + 1-بروموهكسان: \textit{N}-هكسيل فثاليميد؛ ثم يحرر الهيدرازين (أو الحلمهة)
الهكسان-1-أمين. ونيتروجين الفثاليميد المؤلكَل ليس له زوج حر متاح (فهو غير
متمركز على كربونيلين) ولا رابطة \ce{N-H} متبقية: فلا يمكن ألكلته مرة ثانية.
\end{solution}

\begin{solution}{exo:b2:amines:10}
يبرمن ماء البروم الأنيلين في المواضع الثلاثة أورثو وبارا بالنسبة إلى المجموعة
الأمينية: 2،4،6-ثلاثي بروموأنيلين. ثم تستبدل \omterm{def:b2:amines:diazonium}{الديزوة}، يتبعها \ce{H3PO2}، بالمجموعة الأمينية H: فتصير ذرات
البروم الثلاث في المواضع 1،3،5 بعضها بالنسبة إلى بعض.
\end{solution}

\begin{solution}{exo:b2:amines:11}
4-نتروأنيلين + \ce{NaNO2} + \ce{HCl} عند 0--\qty{5}{\degreeCelsius}: كلوريد 4-نتروبنزين ديازونيوم.
ثم الاقتران مع الفينول في محلول ضعيف القاعدية (الفينوكسيد، المنشط بقوة) في
الموضع بارا بالنسبة إلى \ce{O-}: 4-[(4-نتروفينيل)ديازينيل]فينول، \ce{O2N-C6H4-N=N-C6H4-OH}، صبغة
برتقالية.
\end{solution}

\begin{solution}{exo:b2:amines:12}
مجموعة ثلاثي ميثيل الأمونيوم مجموعة مغادرة رديئة وضخمة وتجعل هيدروجينات $\beta$ في
جهة \ce{CH3} أكثر حمضية وأسهل منالًا؛ فتنزع القاعدة هيدروجينًا من الكربون الأقل
استبدالًا: البوت-1-ين، الألكين الأقل استبدالًا (ناتج هوفمان).
\end{solution}

\begin{solution}{pb:b2:amines:1}
\textbf{1.} $\mathrm{{}^-O_3S{-}C_6H_4{-}NH_3^+}$: فمجموعة السلفونيك الحمضية برتنت الأمين.
\textbf{2.} الأيون المزدوج الشحنة قليل الذوبان في الماء؛ والكربونات تنزع البروتون من مجموعة الأمونيوم،
معطية سلفانيلات الصوديوم الذوابة.
\textbf{3.} \ce{NaNO2 + HCl -> HNO2 + NaCl}.
\textbf{4.} لدينا:
\[
  \mathrm{{}^-O_3S{-}C_6H_4{-}NH_2 + HNO_2 + H^+ \to {}^-O_3S{-}C_6H_4{-}N_2^+ + 2\,H_2O} .
\]
\textbf{5.} يتفكك ملح الديازونيوم (إلى الفينول و\ce{N2}) إذا سُخّن؛ ويتفكك حمض النتروز
أيضًا.
\textbf{6.} $5.00/173.19 = \qty{28.9}{mmol}$: \qty{28.9}{mmol} من \ce{NaNO2}، أي \qty{1.99}{g}.
\textbf{7.} حمض النتروز الزائد يتفاعل مع شريك الاقتران؛ ويُكشف عنه بورقة النشاء--اليوديد
ويُتلف باليوريا أو بحمض السلفاميك.
\textbf{8.} حلقته منشطة بقوة بمجموعة ثنائي ميثيل أمينو، وهي مانحة قوية.
\textbf{9.} في الموضع بارا بالنسبة إلى \ce{N(CH3)2}: توجيه أورثو/بارا، والموضعان أورثو تعيقهما
مجموعتا الميثيل.
\textbf{10.} $\mathrm{{}^-O_3S{-}C_6H_4{-}N{=}N{-}C_6H_4{-}N(CH_3)_2}$.
\textbf{11.} في الحمض القوي تُبرتن مجموعة ثنائي ميثيل أمينو: و\ce{-NH(CH3)2+} تثبط الحلقة
فيتوقف الاقتران.
\textbf{12.} يتكوّن الناتج على هيئة صورته الحمضية المبرتنة الذوابة جزئيًا؛ وفي القاعدة يترسب
سلفونات الصوديوم، الأقل ذوبانًا في المحلول المالح.
\textbf{13.} شحنته الموجبة غير متمركزة على الحلقة، والنيتروجين الطرفي مركز محب للإلكترونات
ضعيف؛ فلا تتفاعل إلا الحلقات المنشطة بالمجموعات \ce{-O-} أو \ce{-OH} أو \ce{-NR2}.
\textbf{14.} أصفر برتقالي في القاعدة، وأحمر في الحمض.
\textbf{15.} على أحد نيتروجيني مجموعة الآزو (والصورة المبرتنة يثبّتها الترافق مع
مجموعة ثنائي ميثيل أمينو).
\textbf{16.} نحو $\mathrm pK_a \pm 1$: من 3.1 إلى 4.4 تقريبًا، حول $\mathrm pK_a = 3.47$.
\textbf{17.} تكوّن الحلقتان ومجموعة الآزو نظامًا مترافقًا طويلًا: فالفجوة بين أعلى مستوى
مشغول وأدنى مستوى فارغ من مستويات $\pi$ صغيرة (\cref{prop:b2:huckel:gap}) ويقع الامتصاص
في المرئي.
\textbf{18.} تغيّر البرتنة توزيع الشحنة في نظام $\pi$ وتضيّق الفجوة
أكثر: فينتقل الامتصاص إلى أطوال موجية أطول، واللون من الأصفر إلى الأحمر.
\textbf{19.} \qty{173.19}{g/mol}.
\textbf{20.} \qty{327.33}{g/mol}.
\textbf{21.} \qty{28.87}{mmol}.
\textbf{22.} $0.02887 \times 327.33 = \qty{9.45}{g}$.
\textbf{23.} تفكك ملح الديازونيوم، والاقتران غير التام، وذوبان الناتج في
المحلول الأم وفي مياه الغسل.
\textbf{24.} $0.80 \times 9.45 = \boldsymbol{\qty{7.56}{g}}$ من برتقالي الميثيل.
\end{solution}
